\subsection{Funzioni Test}
\label{subsec:TestRequests}

\subsubsection{run\_test}  
\label{subsubsec:run_test}

\paragraph{Descrizione}  
La funzione \textbf{run\_test} utilizza le domande del dataset indicato per valutare il modello LLM selezionato, restituendo i risultati e le relative statistiche in formato strutturato.


\paragraph{Dettagli dell'endpoint}
\begin{itemize}
    \item \textbf{HTTP Method}: POST;
    \item \textbf{Endpoint}: \texttt{/tests};
    \item \textbf{Response Model}: TestDTO;
    \item \textbf{Dependency Injection}: 
    \begin{itemize}
        \item \textbf{test\_service (TestUseCase)}: la dipendenza viene risolta tramite \textit{Dependency Injector}. Questo permette di ottenere l'istanza del \texttt{TestService} dal \texttt{Container}, che si occupa di tutte le configurazioni e inizializzazioni necessarie per il servizio.
    \end{itemize}
\end{itemize}

\paragraph{Parametri di input}
\begin{itemize}
    \item \textbf{TestDTO}: rappresentazione del test da creare.
\end{itemize}

\paragraph{Funzionamento}
\begin{itemize}
    \item Crea un nuovo test utilizzando le informazioni ricevute;
    \item Se il dataset scelto contiene elementi incompleti restituisce un oggetto \texttt{QaIncompleteListDTO}; 
    \item Restituisce lo stato dell'operazione e nel  caso in cui sia andata a buon fine un oggetto \texttt{TestDTO} che rappresenta il test creato.
\end{itemize}

\subsubsection{get\_test\_page}
\label{subsubsec:get_test_page}

\paragraph{Descrizione}
La funzione \textbf{get\_test\_page} consente di recuperare una pagina di risultati di un test precedentemente eseguito.

\paragraph{Dettagli dell'endpoint}
\begin{itemize}
    \item \textbf{HTTP Method}: GET;
    \item \textbf{Endpoint}: \texttt{/tests/<test\_id>/results};
    \item \textbf{Response Model}: TestPageDTO;
    \item \textbf{Dependency Injection}: 
    \begin{itemize}
        \item \textbf{test\_result\_service (TestResultUseCase)}: la dipendenza viene risolta tramite \textit{Dependency Injector}. Questo permette di ottenere l'istanza del \texttt{TestResultService} dal \texttt{Container}, che si occupa di tutte le configurazioni e inizializzazioni necessarie per il servizio.
    \end{itemize}
\end{itemize}

\paragraph{Parametri di input}
\begin{itemize}
    \item \textbf{id (int)}: ID del test di cui si vuole ottenere una pagina;
    \item \textbf{page (int: 1)}: numero della pagina richiesta.
\end{itemize}

\paragraph{Funzionamento}
\begin{itemize}
    \item Verifica la presenza del test specificato;
    \item Recupera i risultati del test associati alla pagina richiesta;
    \item Restituisce lo stato dell'operazione e nel caso di successo un oggetto \texttt{TestPageDTO} contenente i risultati inclusi nella pagina richiesta.
\end{itemize}

\subsubsection{update\_test}
\label{subsubsec:update_test}

\paragraph{Descrizione}
La funzione \textbf{update\_test} modifica il test indicato.

\paragraph{Dettagli dell'endpoint}
\begin{itemize}
    \item \textbf{HTTP Method}: PUT;
    \item \textbf{Endpoint}: \texttt{/tests/<test\_id>}
    \item \textbf{Response Model}: TestDTO;
    \item \textbf{Dependency Injection}: 
    \begin{itemize}
        \item \textbf{test\_service (TestUseCase)}: la dipendenza viene risolta tramite \textit{Dependency Injector}. Questo permette di ottenere l'istanza del \texttt{TestService} dal \texttt{Container}, che si occupa di tutte le configurazioni e inizializzazioni necessarie per il servizio.
    \end{itemize}

\end{itemize}

\paragraph{Parametri di input}
\begin{itemize}
    \item \textbf{TestDTO}: rappresentazione del test da salvare/modificare.
\end{itemize}

\paragraph{Funzionamento}
\begin{itemize}
    \item Modifica/salvataggio del test indicato;
    \item Restituisce lo stato dell'operazione e se è andata a buon fine un oggetto \texttt{TestDTO} che rappresenta il test salvato. 
\end{itemize}

\subsubsection{delete\_test}
\label{subsubsec:delete_test}

\paragraph{Descrizione}
La funzione \textbf{delete\_test} consente di eliminare un test.

\paragraph{Dettagli dell'endpoint}
\begin{itemize}
    \item \textbf{HTTP Method}: DELETE;
    \item \textbf{Endpoint}: \texttt{/tests/<id\_test>};
    \item \textbf{Response Model}: TestDto;
    \item \textbf{Dependency Injection}: 
    \begin{itemize}
        \item \textbf{test\_service (TestUseCase)}: la dipendenza viene risolta tramite \textit{Dependency Injector}. Questo permette di ottenere l'istanza del \texttt{TestService} dal \texttt{Container}, che si occupa di tutte le configurazioni e inizializzazioni necessarie per il servizio.
    \end{itemize}
\end{itemize}

\paragraph{Parametri di input}
\begin{itemize}
    \item \textbf{id (int)}: ID del test da eliminare.
\end{itemize}

\paragraph{Funzionamento}
\begin{itemize}
    \item Elimina il test indicato;
    \item Restituisce lo stato dell'operazione.
\end{itemize}

\subsubsection{get\_all\_tests}
\label{subsubsec:get_all_tests}

\paragraph{Descrizione}
La funzione \textbf{get\_all\_tests} restituisce l'elenco dei test salvati che contengono la query nel proprio nome o sono stati eseguiti usando un determinato LLM e un determinato dataset.

\paragraph{Dettagli dell'endpoint}
\begin{itemize}
    \item \textbf{HTTP Method}: GET;
    \item \textbf{Endpoint}: \texttt{/tests};
    \item \textbf{Response Model}: TestListDto;
    \item \textbf{Dependency Injection}: 
    \begin{itemize}
        \item \textbf{test\_service (TestUseCase)}: la dipendenza viene risolta tramite \textit{Dependency Injector}. Questo permette di ottenere l'istanza del \texttt{TestService} dal \texttt{Container}, che si occupa di tutte le configurazioni e inizializzazioni necessarie per il servizio.
    \end{itemize}
\end{itemize}

\paragraph{Parametri di input}
\begin{itemize}
    \item \textbf{q (str: "")}: query da utilizzare per filtrare i test salvati;
    \item \textbf{LLM (int)}: ID dell'LLM che deve essere utilizzato nei test salvati da ottenere;
    \item \textbf{dataset (int)}: ID dell'dataset che deve essere utilizzato nei test salvati da ottenere.
\end{itemize}

\paragraph{Funzionamento}
\begin{itemize}
    \item Restituisce lo stato dell'operazione e se è andata a buon fine un oggetto \texttt{TestListDto} la lista dei test salvati.
\end{itemize}

\subsubsection{get\_test\_by\_id}
\label{subsubsec:get_test_by_id}

\paragraph{Descrizione}
La funzione \textbf{get\_test\_by\_id} restituisce le informazioni del test indicato.

\paragraph{Dettagli dell'endpoint}
\begin{itemize}
    \item \textbf{HTTP Method}: GET;
    \item \textbf{Endpoint}: \texttt{/tests/<id\_test>};
    \item \textbf{Response Model}: TestDTO;
    \item \textbf{Dependency Injection}: 
    \begin{itemize}
        \item \textbf{test\_service (TestUseCase)}: la dipendenza viene risolta tramite \textit{Dependency Injector}. Questo permette di ottenere l'istanza del \texttt{TestService} dal \texttt{Container}, che si occupa di tutte le configurazioni e inizializzazioni necessarie per il servizio.
    \end{itemize}
\end{itemize}

\paragraph{Parametri di input}
\begin{itemize}
    \item \textbf{id (int)}: ID del test da ottenere.
\end{itemize}

\paragraph{Funzionamento}
\begin{itemize}
    \item Restituisce lo stato dell'operazione e se è andata a buon fine un oggetto \texttt{TestDTO} che rappresenta il test richiesto.
\end{itemize}

\paragraph{Tracciamento dei requisiti}
\begin{tabularx}{\textwidth}{|X|c|}
    \hline
    \textbf{ID} & \textbf{Funzione} \\
    \hline
        RFO-14.03, RFO-14.06, RFO-14.13, RFO-19.01, RFO-19.02, RFO-19.03  & run\_test \\
    \hline
        RFO-18.03, RFO-18.05, RFF-15.06, RFF-15.08 & get\_test\_page \\
    \hline
        RFF-14.01, RFF-14.04, RFF-15.04, RFF-15.05  & update\_test \\
    \hline
        RFF-14.03 & delete\_test \\
    \hline
        RFF-14.02, RFF-14.06, RFF-20.01 & get\_all\_tests \\
    \hline
        RFO-18.01, RFO-18.02, RFO-18.06, RFO-19.04, RFO-19.05, RFO-19.06, RFO-19.07, RFF-15.01, RFF-15.02, RFF-15.03, RFF-15.07 & get\_test\_by\_id \\
    \hline
\end{tabularx}
